% ЗАКЛЮЧЕНИЕ, без номер.
% По А.3: обобщение на характеристиките и особеностите на решението, предимствата и
% недостатъците му, резултатите от експериментите и тяхната инженерна трактовка, и
% предложения за по-нататъшна работа.
\section*{ЗАКЛЮЧЕНИЕ}
\label{sec:conclusion}
\addcontentsline{toc}{section}{ЗАКЛЮЧЕНИЕ}

В курсовия проект е проектиран и реализиран инструмент за двупосочно инженерство между
изходен код на C++, UML диаграма на класове и релационна схема на данни. Приетият подход
поставя в средата един вътрешен модел, независим както от езика на входа, така и от
формата на изхода, и свежда всяко преобразувание до преход между този модел и една външна
форма. Следствието е, че добавянето на нова нотация не засяга извличането, а промяна в
извличането се отразява едновременно във всички изходни формати.

Предимствата на решението произтичат от избора на компилаторен фронтенд за прочитане на
кода. Понятия като разрешено име, тип след заместване на шаблонните аргументи и
собственост върху обект са достъпни точно защото анализът е семантичен, а не текстов. Това
е и условието разграничението между агрегация и композиция да бъде правено по устройство,
а не по предположение. Второто предимство е в начина на достъп до фронтенда: той се
изпълнява като външен процес, а изводът му се разбира като текст, така че проектът няма
нито една задължителна външна зависимост и се построява само с компилатор за C++20 и
CMake. Ограниченията са в същата посока: инструментът вижда само това, което изводът на
фронтенда съдържа, а при това извличане стандартната библиотека е заместена с декларации,
тоест типове извън заявената папка остават непрозрачни атрибути, а не връзки.

Измереното от четирите експеримента дава следната картина. Всички 20 проверени езикови
конструкции се извличат така, както очакваният модел е записан предварително, а породените
от инструмента заготовки на класове са приети от компилатора. Цената на извличането е под
0.11 s за действителните входове, като изведеното синтактично дърво е около 48 пъти
по-голямо от изходния текст, а полученият от него модел е около 88 пъти по-малък от
дървото, тоест абстракцията е количествено измерима, а не само декларирана. Пълният цикъл
през PlantUML е без загуба и за двата входа, през Mermaid запазва 0.996 от елементите при
заглавните файлове на инструмента и 0.927 при анотирания пример, а през релационната схема
запазва 0.293 при анотирания пример и нула при заглавните файлове. И петте нарочно внесени
разминавания са открити, при нула съобщения, които не отговарят на действително
разминаване.

Три от тези числа заслужават трактовка, а не само отчитане. Първо, разликата между 1.000
и 0.996 не е дефект на модела, а свойство на нотацията: Mermaid отбелязва операция като
абстрактна или като статична и няма трети знак за операция, която е виртуална, но не е
чиста. Именно това е ползата от два излъчвателя вместо един. Второ, нулата през
релационната схема при заглавните файлове на инструмента е вярна: нито един клас на
инструмента не е обозначен за постоянно съхранение, така че въпросът „колко от този модел
може да бъде изразено като таблици“ има отговор нула. Трето, обхождането до модел расте
по-бързо от линейното, защото разрешаването на име търси последователно в списък от всички
класификатори; при 800 класа това вече е най-скъпата от трите съставки на извличането.

Наред с това самите експерименти откриха един действителен дефект, който преглед на кода
не беше открил: безименните параметри не преживяваха пълния цикъл, защото излъчвателят ги
записваше без име, а четецът прочиташе двоеточието обратно като част от типа. Дефектът е
поправен и е покрит с тест. Това е и най-прекият довод в полза на самия подход: щом и
двете посоки са автоматизирани, разминаването става измеримо, а не подозирано, и то важи
както за чужд код, така и за самия инструмент.

За по-нататъшна работа се очертават пет посоки. Първата е замяна на последователното
търсене при разрешаване на имената с речник, което премахва квадратичното поведение,
видимо над 400 класа. Втората е премахване на дублирането на кратността, която днес се пази
и върху връзката, и върху атрибута, поради което една промяна поражда две верни съобщения
вместо едно. Третата е разширяване на покритието на езика към конструкции, останали извън
проверката: приятелски декларации, наследяване от специализация на шаблон и вложени
шаблони. Четвъртата е свързване на проверката за разминаване с процеса на непрекъсната
интеграция, за което инструментът вече е готов, тъй като завършва с ненулев код при
разлика. Петата е допълване на модела с поведенческа информация, което е предпоставка за
извличане, а не само за излъчване, на диаграмите на случаите на употреба и на дейностите.
